% ===== 原文 pp.16–17 =====
% tr-02b：正交出发测地线的三类型、裤子几何（螺线/主间隙/赤道反射）、边界斜交测地线、套索
% 本片含图 6 与图 7（原文同在 p.17），图注照抄 manifest_statements.txt，勿改。

% [原文 p.16]
考虑从 $\beta$ 正交出发的测地线的极大延拓——称这些测地线为\emph{正交出发的测地线}（ortho emanating geodesics）。除正交边界测地线之外，正交出发的测地线有三类：非简单的、无穷长的简单型、以及在第二个端点处与边界斜交的简单型。我们将看到，这些类型可以通过考察一条裤子中的初始段来判别。

考虑由正交边界测地线 $\gamma$ 得到的几何裤子 $\cP$（见图 6）。它的边界是 $\beta$ 以及两条确定的测地线 $\alpha$ 与 $\lambda$（在特殊情形 $(g,n)=(1,1)$ 下 $\alpha=\lambda$）。\emph{螺线}（spiral）是一条无穷的简单测地射线，它聚向一条简单闭测地线。有两条正交出发的测地线是以 $\alpha$ 为聚集集的螺线，另有两条是以 $\lambda$ 为聚集集的螺线。聚向同一条边界的那两条螺线沿该边界朝相反方向缠绕。$\cP$ 的\emph{主间隙}（main gaps）是 $\beta$ 的两条互不相交的子弧，每条子弧的内部含 $\gamma$ 的一个端点，且以螺线的初始点为端点。主间隙将成为与裤子 $\cP$ 相对应的 Cantor 集补集的分量。

由裤子的几何，从主间隙正交出发的测地线或为螺线的端点、或为 $\gamma$、或在 $\cP$ 内自交而非简单、或在第二个端点处与 $\beta$ 斜交而简单。$\beta$ 中主间隙的补集是一对开区间。由裤子的几何，对给定的开区间，所有正交出发的测地线都通过穿越边界 $\alpha$ 或 $\lambda$ 之一而离开这条裤子。对给定的开区间，$\cP$ 中的初始段是简单的，而这些测地线由它们在曲面别处（在某条别的裤子上）的后续行为来分类。连接 $\beta$ 与 $\alpha$ 的正交边界测地线含于其中一个开区间，而从 $\beta$ 到 $\lambda$ 的正交边界测地线含于另一个。一条裤子带有一个赤道反射（equatorial reflection），它稳定每条边界。赤道反射自然地作用在 $\beta$ 的分解上，或交换、或稳定各个元素。

% [Tutor Note: p.16 这段给出的是 \(\beta\) 的一个"四段式"结构，可作为抄写与读图的自检：
%  \(\beta\) 被 2 条主间隙与 2 个开区间交替分成 4 段，每条主间隙的内部恰好含 \(\gamma\) 的一个端点，
%  而主间隙的端点都是螺线的初始点；螺线共 4 条（2 条聚向 \(\alpha\)、2 条聚向 \(\lambda\)），
%  聚向同一底边的两条反向缠绕。图 6 底部圆周即 \(\beta\)，其上加粗的两段弧就是这两条主间隙；
%  读者可用"主间隙数 = 2"与 p.19 的"主间隙长度之和 + 双边界袖口裤子项 = \(\ell_\beta\)"互相校对，
%  若把某条主间隙整段漏算，\(\ell_\beta\) 的等式立刻不成立。]

在上面的讨论中，与一个主间隙相联系的是在裤子内自交的那些正交出发测地线，以及第二次与 $\beta$ 斜交的那些简单测地线——称第二类测地线为\emph{边界斜交}（boundary oblique）测地线。下一个观察是：这些联系可以反转；这些联系给出主间隙与曲面上具有特定行为的测地线之间的双射。

一条边界斜交测地线与 $\beta$ 构成一个歪斜的线框（crooked wire frame），它确定一条裤子，这与正交边界测地线确定一条裤子类似。边界斜交测地线以连续族的形式出现，每个族都收敛到一条正交边界测地线。一个族及其极限确定同一条裤子。族中元素的初始点（即与 $\beta$ 正交的那些点）落在同一个主间隙区间内——这一观察反转了主间隙的各个段与边界斜交测地线的联系。

接下来我们描述如何反转主间隙的各个段与非简单正交出发测地线的联系。这样一条测地线的第一个自交点含于唯一的一条嵌入裤子中。为看出这一点，考虑从 $\beta$ 开始、到该测地线第二次经过其第一个自交点处为止的套索（lasso）子弧。

% [原文 p.17]
\begin{figure}[!htbp]
\centering
\includegraphics[width=0.55\textwidth]{fig06.png}
\figtitle{6}{一个带赤道、主间隙、两条螺线以及一条正交边界测地线 $\gamma$ 的裤子。}
\end{figure}

套索的小邻域/加厚的边界是一条简单闭曲线与 $\cB$ 中一个元素的并。一个几何论证表明：该自交点含于由 $\cB$ 中这个元素所确定的裤子中，而套索的初始点落在该裤子的主间隙内。这一观察反转了主间隙的各个段与非简单正交出发测地线的联系。

% [Tutor Note: 图 7 所画的套索小邻域之所以是关键，是因为它把"自交点"（几何信息）换成了"\(\cB\) 中的一个类"（拓扑信息）：
%  加厚后的边界里那条简单闭曲线负责围住套索圈，而属于 \(\cB\) 的那条弧负责连到底边界，
%  二者合起来才是某条裤子的"线框"。这也解释了为什么非简单测地线的第一个自交点一定是"第一个"才可用：
%  取更靠后的自交点，套索会先绕出当前裤子，加厚边界就不再是"一条简单闭曲线 + 一个 \(\cB\) 元素"了。]

\begin{figure}[!htbp]
\centering
\includegraphics[width=0.42\textwidth]{fig07.png}
\figtitle{7}{一条套索（lasso）的小邻域。}
\end{figure}

我们回忆：带内部交点的非简单性，在测地线空间上是开条件，在正交出发测地线空间上也是开条件。通过考虑曲面的加倍，“与所有边界交点都正交的简单性”是测地线空间上的闭条件。而“与边界正交且只相交一次的简单性”之集是一个 Cantor 集。正交出发测地线的分类至此完成。

% [Tutor Note: p.17 末这两句是全篇长度恒等式的枢纽：定理 8 的左端 \(L_1\) 等于 \(\beta_1\) 上 Cantor 集
%  的补区间（即主间隙）长度之和，而"补区间之和 = 全长"这件事恰恰依赖此处的开/闭条件分类——
%  开条件保证四类点集互不相交且覆盖，闭条件保证 Cantor 集（正交单交的简单点集）确实闭。
%  注意定理 7 的两条结论并不矛盾：测地线空间中的集合 \(S\) 的 Hausdorff 维数为 1（因为含曲面内部的
%  无穷长测地线），而 \(S\) 在边界 \(\beta\) 上的截口集维数与测度都是 0——求长时只用后者。]
